% SGA 6, Exposición X: fragmento español alineado con la fuente.
% Alcance: sección 6 completa, «Apéndice: Cálculo de los determinantes
% de los haces localmente libres».
% Fuente francesa: sga6_fr_r1/src/fr.tex, líneas 13296--13494.
% SHA-256 de la fuente francesa completa:
% 0D2158B9FE3B466309A92165727A0B191422E05FBAA0DACD5806E8322C54BA1A.

\section*{6. Apéndice: Cálculo de los determinantes de los haces localmente libres}

\subsection*{6.1}
Consideremos primero un topos anillado cualquiera
\(X=(\mathcal C,\mathcal O)\). Para toda sección positiva
\(n\in\Gamma(\mathbb Z_X)\) del haz constante \(\mathbb Z\) y para todo
\(\mathcal O\)-módulo \(M\in\Mod_{\mathcal O}(X)\), está claro cómo definir
los \(\mathcal O\)-módulos
\begin{equation}\tag{6.1.1}
 M^n,\qquad \bigotimes^n M,\qquad \bigwedge^n M,\quad\text{etc.}
\end{equation}
Los funtores \((\ )^n\), \(\bigotimes\limits^n\) y
\(\bigwedge\limits^n\) transforman módulos libres (resp. localmente libres)
en módulos libres (resp. localmente libres).

En general, consideremos secciones \(n_1,\ldots,n_r\) de \(\mathbb Z_X\) y
un multifuntor (covariante, para fijar ideas)
\[
 T:\Mod^{n_1}(\mathcal O)\times\cdots\times
 \Mod^{n_r}(\mathcal O)\longrightarrow\Mod(\mathcal O)
\]
de \(r\) variables, donde \(\Mod_{\mathcal O}^{n}(X)\) designa la categoría
fibrada sobre \(X\) de los \(\mathcal O\)-módulos localmente libres de rango
\(n\). Sea
\[
 M=T(\mathcal O^{n_1},\ldots,\mathcal O^{n_r})
\]
y
\[
 G=\operatorname{G\ell}(n_1)\times\cdots\times\operatorname{G\ell}(n_r),
 \quad\text{donde }\operatorname{G\ell}(n_i)=\underline{\Aut}_{\mathcal O}(\mathcal O^{n_i}).
\]
Se sigue inmediatamente que \(T\) define una representación
\begin{equation}\tag{6.1.2}
 \alpha_T:G\longrightarrow\underline{\Aut}_{\mathcal O}(M).
\end{equation}
Está claro que, si \(M\) es a su vez localmente libre de rango \(m\), para
todo funtor compuesto
\[
 T'T,\ \text{donde } T':\Mod^{m}(\mathcal O)\longrightarrow\Mod(\mathcal O),
\]
se tiene
\begin{equation}\sisi{\tag{6.1.3}}{} % Canon SGA-SRC-001341
 \alpha_{T'T}=\alpha_{T'}\,\alpha_T.
\end{equation}

Recíprocamente, dada la representación
\(\alpha=\alpha_T:G\to\underline{\Aut}_{\mathcal O}(M)\), se recupera
\(T(E_1,\ldots,E_r)\) (para toda sucesión \(E_1,\ldots,E_r\) de
\(\mathcal O\)-módulos localmente libres de rangos \(n_1,\ldots,n_r\))
pegando el \(\mathcal O\)-módulo libre
\(T(\mathcal O^{n_1},\ldots,\mathcal O^{n_r})\) sobre cartas locales
trivializantes mediante la representación \(\alpha\). Más precisamente, a
la sucesión \(E_1,\ldots,E_r\) le corresponde un \(G\)-torsor derecho
\(P=P_{E_1,\ldots,E_r}\) tal que se tiene un isomorfismo canónico
\begin{equation}\tag{6.1.3}
 E_i\xrightarrow{\sim}P\overset{G}{\wedge}\mathcal O^{n_i}\qquad(i=1,\ldots,r),
\end{equation}
donde \(\overset{G}{\wedge}\) designa el producto contraído, y \(G\) actúa
sobre \(\mathcal O^{n_i}\) mediante la \(i\)-ésima proyección
\(G\to\Gl(n_i)\); entonces toda representación
\[
 \alpha:G\longrightarrow\underline{\Aut}_{\mathcal O}(M)
\]
hace actuar a \(G\) sobre \(\sisi{\mathcal O^m}{M}\) mediante % Canon SGA-SRC-001342
\(\mathcal O\)-automorfismos y define así un multifuntor \(T=T_\alpha\) tal
que
\begin{equation}\tag{6.1.4}
 T_\alpha(E_1,\ldots,E_r)=P_{E_1,\ldots,E_r}\overset{G}{\wedge} M.
\end{equation}
Además, para toda representación compuesta
\(\alpha=\alpha_1\circ\alpha_2\), se obtiene
\begin{equation}\tag{6.1.5}
 T_\alpha=T_{\alpha_1}\circ T_{\alpha_2}.
\end{equation}

\subsection*{6.2}
Para todo \(\mathcal O\)-módulo localmente libre \(E\) de rango \(n\),
designaremos por \(\det(E)\) el \(\mathcal O\)-módulo invertible
\(\bigwedge\limits^n E\). Conservando las notaciones de 6.1, el funtor
\(\det(T(E_1,\ldots,E_r))\) queda determinado, salvo isomorfismo canónico,
por la representación correspondiente
\begin{equation}\tag{6.2.1}
 \det\circ\alpha_T:G\longrightarrow\Gm=\Gl(1).
\end{equation}

Supongamos ahora que \(X\) es el topos anillado de los haces sobre el sitio
anillado \((\Sch/S)\), provisto de la topología étale. Entonces se puede
calcular \(\det(T(E_1,\ldots,E_r))\), o, lo que es lo mismo, la
representación \(\det\circ\alpha_T\), con ayuda de la

\textbf{Proposición 6.2.2.} Para toda sección positiva \(n\) de
\(\mathbb Z\), el haz étale de grupos \(\Der(\Gl(n))\) de los conmutadores
de \(\Gl(n)\sisi{)}{}\) es representable por % Canon SGA-SRC-001343
\(\operatorname{Sl}(n)=\Ker(\det:\Gl(n)\to\Gm)\).

\textbf{Corolario 6.2.3.} Toda representación
\begin{equation}\tag{6.2.4}
 \alpha:G=\Gl(n_1)\times\cdots\times\Gl(n_r)\longrightarrow\Gm
\end{equation}
es de la forma
\begin{equation}\tag{6.2.5}
 \alpha(Y)(s)=(\det(s_1))^{m_1}\cdots(\det(s_r))^{m_r},
\end{equation}
donde \(Y\in\operatorname{Ob}(\Sch/S)\),
\(s=(s_1,\ldots,s_r)\in G(Y)\), y \(m_1,\ldots,m_r\) son secciones de
\(\mathbb Z\) determinadas unívocamente por \(\alpha\).

\textbf{Corolario 6.2.6.} Con las notaciones de 6.1, se tiene
\begin{equation}\tag{6.2.6}
 \det\bigl(T_\alpha(E_1,\ldots,E_r)\bigr)
 =
 \det(E_1)^{\otimes m_1}\otimes\cdots\otimes
 \det(E_r)^{\otimes m_r},
\end{equation}
donde \(m_1,\ldots,m_r\) son secciones de \(\mathbb Z\) determinadas
unívocamente por \(\alpha\).

\textbf{Observación 6.2.7.} Se pueden determinar los exponentes
\(m_1,\ldots,m_r\) en (6.2.5) y (6.2.6) sustituyendo sucesivamente en
(6.2.5)
\(s_i=\lambda.\ \operatorname{id}.(\lambda\in\Gm(Y))\),
\(s_j=\operatorname{id}_{\Gl(n_j)(Y)}\), para \(j\ne i\)
\((i,j=1,\ldots,r)\), de donde
\[
 \det(s_i)=\lambda^{n_i}\quad\text{y}\quad
 \alpha(Y)(s)=\lambda^{n_im_i}.
\]
Por ejemplo, se deduce así de (6.2.6)
\begin{align}
 \det(\sisi{E_i}{E_1}\otimes E_2)&=(\det(E_1))^{\otimes n_2}\otimes(\det(E_2))^{\otimes n_1}, \tag{6.2.8}\\ % Canon SGA-SRC-001344
 \det(E_1\sisi{\otimes}{\oplus} E_2)&=(\det(E_1))\otimes(\det(E_2)) \tag{6.2.9} % Canon SGA-SRC-001345
\end{align}
y
\begin{equation}\tag{6.2.10}
 \det\bigl(\bigwedge\limits^m E_1\bigr)=(\det(E_1))^{\binom{n_1-1}{m-1}}
\end{equation}
etc.

\subsection*{6.3}
Mostremos primero cómo se deduce 6.2.3 de 6.2.2. Para ello basta
evidentemente con demostrar (6.2.5) para los
\(s_{(i)}=(1,\ldots,s_i,\ldots,1)\) \((i=1,\ldots,r)\); por tanto, se puede
suponer \(r=1\), \(G=\operatorname{G\ell}(n_1)\). Según 6.2.2, \(\alpha\)
se factoriza a través de
\(\operatorname{G\ell}(n_1)/\operatorname{Sl}(n_1)\), es decir, existe un
diagrama conmutativo
\[
\begin{tikzcd}[column sep=large,row sep=large]
\operatorname{G\ell}(n_1) \arrow[r,"\alpha"] \arrow[d,"j"'] & \Gm\\
\operatorname{G\ell}(n_1)/\operatorname{Sl}(n_1) \arrow[ur,"\beta"'] &
\end{tikzcd}
\]
donde \(j\) es el epimorfismo canónico y \(\beta\) una representación de
\(\operatorname{G\ell}(n_1)/\operatorname{Sl}(n_1)\) en \(\Gm\). Pero el
homomorfismo
\[
 \det:\operatorname{G\ell}(\sisi{n_i}{n_1})\longrightarrow\Gm % Canon SGA-SRC-001346
\]
posee una sección, de modo que el monomorfismo
\[
 \operatorname{G\ell}(\sisi{n_i}{n_1})/\operatorname{S\ell}(\sisi{n_i}{n_1})\longrightarrow\Gm
\]
inducido por det es un isomorfismo. Por otra parte, se sabe (SGA 3 VIII 1.3)
que toda representación de \(\Gm\) en sí mismo es de la forma
\[
 s\rightsquigarrow s^m,
\]
donde \(m\) es una sección de \(\mathbb Z\), de donde resulta la afirmación.

El corolario 6.2.6 es una simple reformulación de 6.2.3. Obsérvese que 6.2.3
se deduce también inmediatamente de SGA 3 XXII 6.2.1, pero queremos dar aquí
una demostración elemental mediante 6.2.2\sisi{)}{}. % Canon SGA-SRC-001347

\subsection*{6.4. Demostración de 6.2.2}
Basta con demostrar que, cualesquiera que sean
\(Y\in\operatorname{Ob}(\Sch/S)\), \(y\in Y\) y
\(s\in\operatorname{S\ell}(n)(Y)\), se puede encontrar un «~entorno
étale~» \(Y'\to Y\) de \(y\) tal que la restricción de \(s\) a \(Y'\) sea
un producto de conmutadores de \(\operatorname{G\ell}(n)(Y')\). Al ser la
cuestión local, se puede suponer que \(Y\) es afín, de anillo \(A\), y que
\(n\) es constante e igual a \(n\) sobre \(Y\). Sean \(\underline p\) el
ideal primo de \(y\) y \((s_{ij})(i,j=1,\ldots,n)\) la matriz de \(s\).
Para todo \(\lambda\in A\) y \(h, k=1,\ldots,n\); \(h\ne k\), designaremos
por \(c_{h,k,\lambda}\) la matriz \((c_{ij})\)
\((i,j=1,\ldots,n)\) tal que \(c_{hk}=\lambda\) y
\(c_{ij}=\delta_{ij}\) para \((i,j)\ne(h,k)\) \((\delta_{ij}=\) símbolo de
Kronecker). Para toda matriz \(a=(a_{ij})(i,j=1,\ldots,n)\), se obtiene
\(c_{h,k,\lambda}a\) añadiendo \sisi{a la \(k\)-ésima columna de \(a\)
\(\lambda\) veces la \(h\)-ésima}{a la \(h\)-ésima fila de \(a\)
\(\lambda\) veces la \(k\)-ésima}. Del mismo modo, se obtiene la matriz % Canon SGA-SRC-001349
\(ac_{h,k,\lambda}\) añadiendo \sisi{a la \(h\)-ésima fila de \(a\)
\(\lambda\) veces la \(k\)-ésima}{a la \(k\)-ésima columna de \(a\)
\(\lambda\) veces la \(h\)-ésima}. En particular, se tiene
\[
 \det(c_{h,k,\lambda})=1,
\]
y se ve inmediatamente que
\begin{equation}\tag{\sisi{6.3.1}{6.4.1}} % Canon SGA-SRC-001348
 c_{h,k,\lambda}c_{h,k,\mu}=c_{h,k,\lambda+\mu}
\end{equation}
para todo \(\lambda,\mu\in A\), de modo que se obtiene una representación
\[
 c_{h,k}:G_a\longrightarrow\operatorname{S\ell}(n)
 \qquad(h,k=1,\ldots,n\ ;\ h\ne k).
\]
\sisi{Salvo indicación expresa en contrario, en lo que sigue supondremos
invertible el \(\lambda\) que interviene en \(c_{h,k,\lambda}\)
\((\lambda\in A\ ;\ h,k=1,\ldots,n\ ;\ h\ne k)\).}{No es necesaria ninguna hipótesis de invertibilidad sobre el parámetro
\(\lambda\) en el primer paso de la reducción matricial; nos reduciremos a
parámetros invertibles antes del segundo paso, dedicado a los conmutadores.}

La demostración se hace entonces en dos pasos. Mostraremos primero que \(s\)
es un producto de matrices \(c_{h,k,\lambda}\) en un entorno zariskiano
suficientemente pequeño de \(y\). Después mostraremos que las matrices
\(c_{h,k,\lambda}\) son productos de conmutadores de
\(\operatorname{G\ell}(n)\) en un entorno étale suficientemente pequeño de
\(y\).

Para el primer paso, se trata de transformar \(s\) en la matriz unidad
\(\mathbf 1=(\delta_{ij})\) en un entorno zariskiano de \(y\), multiplicando \(s\)
a izquierda y a derecha por matrices adecuadas \(c_{h,k,\lambda}\).
Mostremos primero que de este modo se puede hacer \(s_{11}=1\). Si
\(s_{11}\ne1\) y \(s_{21}\in\underline p\), buscamos una matriz
\(c=\sisi{c_{h,1,\lambda}}{c_{2,h,\lambda}}\) tal que el producto \(\sisi{s'=sc}{s'=cs}=(s'_{ij})\) tenga
\(s'_{21}\notin\mathfrak p\). Puesto que
\(\det(s)=1\notin\mathfrak p\), existe al menos un índice \(i\) tal que
\(s_{i1}\notin\underline p\). Entonces
\(s_{21}+s_{i1}\notin\underline p\), y \(\sisi{c_{i,1,1}}{c_{2,i,1}}\) resuelve el problema.
Se puede, pues, suponer \(s_{21}\notin\underline p\) y, tras localizar, se
puede suponer que \(s_{21}\) es invertible. Si
\(s_{11}-1\in\underline p\), entonces en
\(s'=(s'_{ij})=\sisi{sc_{1,2,1}}{c_{1,2,1}s}\) se tiene
\(s'_{11}-1=(s_{11}+s_{21})-1\notin\underline p\). Tras localizar, se puede
entonces suponer que \(s_{11}-1\) es invertible. Así,
\(\lambda=(1-s_{11})/s_{21}\) es invertible y se obtiene \(s'_{11}=1\) en
\(s'=(s'_{ij})=\sisi{sc_{1,2,\lambda}}{c_{1,2,\lambda}s}\).

Suponiendo \(s_{11}=1\), se puede entonces dar a \(s\) la forma en que
\[
 s_{11}=1,\qquad s_{1i}=s_{i1}=0\quad(i=2,\ldots,n)
\]
en un entorno zariskiano de \(y\), multiplicando \(s\) a izquierda (resp. a
derecha) por matrices \(c_{i,1,\lambda}\) (resp.
\(\sisi{c_{i,1,\lambda}}{c_{1,i,\lambda}}\)). A continuación se hace \(s_{22}=1\), etc. Por
recurrencia, se puede entonces transformar \(s\), en un entorno zariskiano
\(U\) de \(y\), en una matriz \(s'=(s'_{ij})\) tal que
\(s'_{ij}=\delta_{ij}\) para \(i<n\) o \(j<n\). Pero, por la propia
construcción, se tiene entonces \(s'=c_1sc_2\) sobre \(\mathcal U\), donde
\(c_1\) y \(c_2\) son productos de matrices \(c_{h,k,\lambda}\). Por tanto,
se tiene
\[
 \det(s')=\det(s)=1,
\]
de modo que \(s'_{nn}=1\), \(s'=\mathbf 1\), y hemos terminado.

\sisi{}{Tras restringir a un entorno de \(y\), un parámetro
\(a\notin\mathfrak p\) es invertible. Si \(a\in\mathfrak p\), entonces
\(a=1+(a-1)\), y \(1\) y \(a-1\) son invertibles en un entorno de \(y\);
por (6.4.1),
\[
 c_{h,k,a}=c_{h,k,1}c_{h,k,a-1}.
\]
} % Canon SGA-SRC-001349

Suponemos, pues, que \(s\) ya es sobre \(Y\) un producto de matrices
\(c_{h,k,\lambda}\).

Queda por ver que estas son productos de conmutadores en un entorno étale de
\(y\). Ante todo, se ve que dos matrices cualesquiera
\(c_{h,k,\lambda}\) y \(c_{h,k,\mu}\) son conjugadas; en otras palabras,
existe una matriz invertible \(a=(a_{ij})\)
\((i,j=1,\ldots,n)\) tal que
\[
 ac_{h,k,\lambda}=c_{h,k,\mu}a.
\]
En efecto, la matriz \((a_{ij})\) con \(a_{hh}=1/\sisi{\mu}{\lambda}\),
\(a_{kk}=1/\sisi{\lambda}{\mu}\), y en lo restante \(a_{ij}=\delta_{ij}\), resuelve el % Canon SGA-SRC-001350
problema (recuérdese que, por hipótesis, \(\lambda\) y \(\mu\) son
invertibles y \(h\ne k\)). Por tanto, para un homomorfismo cualquiera
\(\beta\) de \(\operatorname{G\ell}(n)(U)\) en un haz de grupos abelianos
\(B\), se tiene
\[
 \beta(c_{h,k,\lambda})=\beta(c_{h,k,\mu}).
\]
Si \(\lambda+\mu\) es invertible, se deduce de esto y de (\sisi{6.3.1}{6.4.1}) que
\(\beta(c_{h,k,\lambda})=0\) para \sisi{todo \(h,\beta\)}{todos \(h,k\)}, de modo que % Canon SGA-SRC-001351
\(c_{h,k,\lambda}\) es un producto de conmutadores. Si el cuerpo residual
\(k\) de \(y\) no es el cuerpo primo
\(\underline F_2=\{0,1\}\), entonces se pueden encontrar, para todo
\(\lambda\in A\) invertible, un entorno zariskiano \(U=\Spec(B)\) de \(y\)
y un elemento invertible \(\mu\in B\) tal que \(\lambda+\mu\in B\) sea
invertible. De ello se sigue que \(c_{h,k,\lambda}\) es un producto de
conmutadores sobre \(U\).

Si \(k=\underline F_2\), entonces \(V=\Spec(A[T]/(T^2+T+1))\) es un entorno
étale de \(y\), cuyo cuerpo residual es distinto de \(\underline F_2\).
Hemos terminado.

\subsection*{6.5. Observaciones}
\textbf{6.5.1.} Cuando ningún cuerpo residual de \(S\) es isomorfo a
\(\underline F_2\), la demostración anterior muestra que 6.2.2
sigue siendo válida para el topos anillado de los haces zariskianos sobre
\(S\). En cambio, ya para \(S=\Spec(\mathbb F_2)\),
\(\operatorname{S\ell}(2,\mathbb F_2)=\operatorname{G\ell}(2,\mathbb F_2)\)
es distinto del grupo de los conmutadores de
\(\operatorname{G\ell}(2,\underline F_2)\). El orden de este último es \(3\),
mientras que el orden de \(\operatorname{S\ell}(2,\mathbb F_2)\) es \(6\)
(cf. [1]).

\textbf{6.5.2.} En un topos localmente anillado cualquiera \(X\), no es
cierto en general que toda representación \(\Gm\to\Gm\) esté definida por
una potencia. Por ejemplo, en el topos \((\Ens)\) anillado por el cuerpo
\(\mathbb Q\), la representación \(\Gm\to\Gm\) definida por
\[
 r\rightsquigarrow|r|\qquad(r\in\mathbb Q^*)
\]
no es una potencia. En particular, 6.2.3 no sigue siendo válido para
cualquier topos localmente anillado \(X\) y para cualquier representación
\(\alpha:G\to\Gm\).

\textbf{6.5.3.} En cambio, 6.2.3 sigue siendo válido para las
representaciones asociadas a las fórmulas (6.2.8), (6.2.9) y (6.2.10) sobre
un topos anillado cualquiera \(X=(\mathcal C,\mathcal O)\), porque estas
representaciones son «~algebraicas~» en el sentido siguiente:

Sean \(e\) el objeto final de \(X\), \(S=\Spec(\mathcal O(e))\), y
consideremos el funtor
\[
 u:\mathcal C\longrightarrow(\Sch/S)
\]
que asocia a cada \(Y\in\operatorname{Ob}(\mathcal C)\) el esquema afín
\(\Spec(\mathcal O(Y))\) sobre \(S\). Entonces, para todo entero positivo
\(n\), el prehaz \(\operatorname{G\ell}(n)_X\) sobre \(\mathcal C\) es el
compuesto con \(u\) del prehaz correspondiente
\(\operatorname{G\ell}(n)_S\) sobre \((\Sch/S)\). Se dice que una
representación sobre \(X\)
\begin{equation}\tag{6.5.4}
 \alpha:\operatorname{G\ell}(n)\longrightarrow\Gm
\end{equation}
es \emph{algebraica} si proviene mediante \(u\) de una representación
\begin{equation}\tag{6.5.5}
 \alpha_S:\operatorname{G\ell}(n)_S\longrightarrow G_{m,S}
\end{equation}
de prehaces en grupos sobre \((\Sch/S)\), es decir, si se tiene
\[
 \alpha(s)=\alpha_S(s)
\]
para todo \(Y\in\operatorname{Ob}(\mathcal C)\),
\(s\in\operatorname{G\ell}(n)(Y)=\operatorname{G\ell}(n)_S(u(Y))\). La
representación \(\alpha\) de (6.5.4) está, por supuesto, completamente
determinada por \(\alpha_S\). Se sigue inmediatamente que (6.2.3) continúa
siendo válido en un topos anillado cualquiera cuando \(\alpha\) es
algebraica. Por tanto, en particular, las fórmulas (6.2.8), (6.2.9) y
(6.2.10) continúan siendo válidas sobre un topos anillado cualquiera.
